% New scoped proof audit, 2026-10-09. No frozen source is changed.
\subsection{The current eighth-prefix weight and the moving-target product estimate}
\label{sec:current-eighth-distribution-interface-audit-20261009}

\paragraph{Finding and scope.}
The delicate implication is not that an eighth-prefix integer has a
factorization. It is that the \emph{same signed finite coefficient}
has a finite labelled convolution expansion whose slot partition is
fixed throughout each box, and whose factors satisfy the actual
moving-residue hypotheses of L23 Proposition 5.2. I re-read the
versioned primary statement and its proof, together with the present
construction, rather than infer this implication from an earlier
audit's status. The implication below is valid for the stated fixed
geometric product families and buffered finite policies. I found no
new unmatched distribution hypothesis for mask138. An important
qualification is proved explicitly below: the carrier level
$\Theta=299/500$ is not itself a licence for every high-rank atom.

The ledger inspected is
\texttt{eighth-history-curvature-transfer-assembly-20261009.json}:
87 selected eighth domains and $K=67617373/10^7$. This review does
not recompute their polytopes, floors, density total or Chen lower
functional. It is not another numerical certificate of that $K$,
nor a proof of the external spectral theorem. Its conclusion is the
analytic applicability of the stated distribution and positive-mass
steps, conditional on the actual finite density bounds with their
already reserved losses. No file of the preceding ledger is changed.

\paragraph{1. Precisely the source theorem.}
Use $T$ for the even target, $x$ for a dyadic outer-product scale,
$U$ for the selected individual Heath--Brown variable, and
$\sigma=7/32$. This avoids the three different uses of $N$ in the
source and project. The primary text is L23,
\href{https://arxiv.org/abs/2309.08522v1}{arXiv:2309.08522v1},
saved as \texttt{external-source-audit/l23-v1-joint-current.txt}.
Definition2.1 is at lines1239--1259; Proposition5.2 is at
2411--2452 and its proof starts at2453. For $UV\asymp x$, it requires
\begin{align}
 Q_1&<Ux^{-\epsilon},& U^2Q_2Q_3^2&<x^{1-8\epsilon},\notag\\
 |a|^\sigma UQ_1Q_2^5Q_3^2&<x^{2-15\epsilon},&
 U^{1+2\sigma}Q_1Q_2^{5-2\sigma}Q_3^{2+2\sigma}
   &<x^{2-15\epsilon}.                         \label{eq:current8-source-sizes}
\end{align}
Its two outer sequences are one-bounded; the $U$ sequence is rough
above $Z=x^{1/(\log\log x)^3}$ and satisfies the full
auxiliary-coprimality Siegel--Walfisz condition. The three modulus
sequences are arbitrary one-bounded complex sequences supported on
$(q_i,a)=1$. The conclusion is uniform over $0<|a|<x^{1+\epsilon}$.
It is an absolute estimate of the signed triple modulus sum, not a
termwise progression estimate or an absolute sum over all moduli.

Target dependence of the three modulus sequences is allowed: the
proof uses their uniform bounds and coprimality, and the product
structure. In particular its factorization of $q_2,q_3$ preserves
independence from $q_1$ (primary lines2453--2520). Dependence on $T$
does not change this. Dependence coupling the three factors to the
individual outer summation variables would change it and is not
licensed. Uniform SW constants for the selected outer sequence are
essential; arbitrary target-dependent bounded outer coefficients
are not covered.

Proposition6.1 does not literally state an eighth-mask theorem.
The project extension uses the elementary construction behind
Lemma6.3 (primary lines2867--2935), proved for an arbitrary integer
selected product, together with the full and remaining-square
conditions below. The external theorem being used is versioned;
this note makes no claim about verification of a later journal
edition or about a stronger unnamed distribution theorem.

\paragraph{2. Two different third cutoffs must not be conflated.}
The actual rule in \texttt{extended-prefix-upper.tex}, paragraphs
``Margins at the larger first stop'' and ``Integer support, weights,
and distribution'', is as follows. Inside the carrier
$\mathcal D^+(T^{\Theta-2e})$, the retained first label lies below
$\rho-4e$, $\rho=19/96$. A third coarse label has the cap
\[
 \lambda_B=\min\{\Theta-2e,\Gamma_B-2e\}.
\]
This cap is imposed whenever the third prime is above the presieve
cutoff $w=T^\omega$. Below $w$ the raw carrier need not have that
extra cap. On the lower first-prime region $\Gamma_B$ is the robust
gated three-prefix envelope; elsewhere it is the single-prime
$v(t_1)$ fallback. Both are at most $\Lambda=153/256$.

The numerical split $t_3=1/100$ in the old density quadrature, and
the same activation restriction retained for the new eighth
patches, are \emph{not} an instruction to omit the base third cap
on $w\le p_3<T^{1/100}$. The density tail below $1/100$ bounds
that capped region's integral. The integer support is the one just
stated, not the truncated numerical integration domain.

Consequently every compiled rough atom has one of two forms:
\begin{enumerate}
\item It has at least three rough slots. Its third prime is at least
the rough cutoff, so the base third cap is already imposed. Any
later active selected cap, including mask138, is below its tent
$v(s)\le\Lambda$, with a fixed buffer. The atom has an actual
licensed high-rank cap and complete well support at that cap.
\item It has zero, one or two rough slots. Its rough product is at
most $T^{2\rho+o(1)}$; including the small beta factor still leaves
a fixed exponent below $1/2$. It is treated by classical/generalized
BV, not by pretending that the carrier $\Theta$ passes Lemma6.3.
\end{enumerate}
Use consistent half-open endpoint conventions for the presieve and
coarse labels. Assigning a prime exactly at the rough threshold to
the capped side preserves the same parity proof; alternatively its
reciprocal-prime density is $O(1/w)$ and can be paid before passage
to the limit. No boundary prime is allowed to become an unlicensed
high-rank atom by a strict/non-strict mismatch.

This distinction matters. For an unbuffered generic level
$\vartheta=\Theta>\Lambda$, the critical interval in Lemma6.3
already has empty width in its residue-dependent constraint:
\[
 \frac{2-2\vartheta-\sigma}{3}-(2\vartheta-1)
   =\frac{5-\sigma-8\vartheta}{3}<0.
\]
Thus $\Theta<3/5$ alone is \emph{not} a proof of the high-rank
factorization. The actual rough-rank split above supplies the
missing reason. A coarse global support bound
$T^{\Theta+\ell+O(M\tau)}<T^{3/5}$ remains a valid convenient
bound if $\ell+O(M\tau)<1/500$; it is not substituted for the
factorization argument.

\paragraph{3. The actual mask138 conditions and all descendants.}
Let $t_1>\cdots>t_8>0$, $S_j=\sum_{i\le j}t_i$,
$H_j=S_{j-1}+3t_j$, and $b=t_2+t_4+t_8$.
Besides the eight initial full squares, a new cap $m_8$ for
$I=\{2,4,8\}$ must satisfy the following five tests:
\begin{gather*}
 b+2t_1<m_8,\quad b+t_1+2t_3<m_8,\quad
 b+t_1+t_3+2t_5<m_8,\\
 b+t_1+t_3+t_5+2t_6<m_8,\quad
 b+t_1+t_3+t_5+t_6+2t_7<m_8,\qquad m_8<v(b).
\end{gather*}
These are the complement slots $1,3,5,6,7$, not merely the first
remaining square. They appear explicitly in
\texttt{eighth-unequal-parent-interface-20261009.tex}, (U5), and in
the latest first-failure physical-row proof.

For every later slot $j>8$, the selected-remaining test equals
$S_{j-1}+2t_j$. An admitted odd cube bounds that square. At even
$j>8$, ordering gives
$S_{j-1}+2t_j\le S_{j-2}+3t_{j-1}$, the preceding admitted odd
cube. Thus all later tests hold at the frozen cap. This is exactly
the complete support required by the selected-product construction.
The eighth node stays admitted; its ordinary upper residual has
odd first exits at full ranks9,11,\ldots. No even first exclusion
is introduced.

The latest29 ``first-failure'' cells use $B=A(t_1)\le L_3$, not
the older $D(t_2+t_3)$ minorant, for all actual histories.
Their failure is relative to that older sufficient condition,
not an actual earlier tree exit. The main ordered-envelope lemma
and its robust coarse-cap rule supply, on every common compact
trim, one first-three-label function $\lambda_3$ with
$H_3,H_5,H_7<\lambda_3$. It is not chosen independently for an
eighth child. With the updated sixth policy,
\[
 \lambda_3\le\lambda_6^{\rm cur}\le R_6<m_8,
 \qquad
 \lambda_6^{\rm cur}-S_8>2t_7-t_8>t_8.
\]
The retained range $m_8-S_8<3t_8$ gives both ordinary arguments
in $(1,3)$. The signed density identity is the difference of the
two actual upper residuals. Its conservative density is
\[
 \frac{4(m_8-R_6)}{(R_6-S_8)(m_8-S_8)\prod_{i=1}^8t_i}.
\]
This uses ordinary raw-density \emph{equality}; it is not obtained
by subtracting two one-sided sieve estimates. The raw-density
source qualification of the existing Suzuki interface is retained,
not re-proved here. Moment or curvature refinements estimate this
same density and do not change any of these supports.

\paragraph{4. From those gates to a fixed convolution partition.}
At a local scale use a slightly larger buffered support exponent
$\vartheta$ and $D=x^{\vartheta-2\delta}$, with
$1/2\le\vartheta<5/8$. For the selected endpoint product $b_0$, put
\[
 v_0=2\vartheta-1,\qquad
 u_0=\min\left\{\frac{2-(2+2\sigma)\vartheta}{3-4\sigma},
                   \frac{2-2\vartheta-\alpha\sigma}{3}\right\}.
\]
Here $T\le x^\alpha$ and $\alpha=1+O(\kappa)$; it is not set
equal to one on the smaller local scales. Strict tent/gate buffers
give $x^{v_0}\le b_0\le x^{u_0}$ and all remaining squares
at most $D$. These statements survive the small coordinate
change from $\log_T$ to $\log_x$.

For $A_0$ in the source range let $C=D/(A_0b_0)$. If $C\ge1$,
put selected slots in group two and greedily assign the remaining
ordered slots into capacities $A_0,C$. If a next slot $q$ fits
neither capacity, then
$q^2>(A_0/a)(C/c)=D/(ab_0c)$, contradicting its remaining-square
condition. Hence $a\le A_0,c\le C,b=b_0$. Direct substitution gives
\[
 A_0^2bc^2\le x^{1-4\delta},\quad
 T^\sigma A_0ab^5c^2\le x^{2-4\delta},\quad
 (A_0c/b)^{2\sigma}A_0ab^5c^2
       \le x^{2-(4+4\sigma)\delta}.
\]
If $C<1$, keep neither the selected assignment nor an imaginary
capacity below one. Reassign \emph{all} slots into capacities
$A_0,B=D/A_0$, using the full-square conditions and the same
two-capacity contradiction. Here $B<b_0\le x^{u_0}$ and
$B\ge1$. With $c=1$, the first monomial is at most
$A_0D\le x^{23/24-5\delta/2}$; the other two are bounded by
$x^{\alpha\sigma}D^2B^3$ and
$D^{2+2\sigma}B^{3-4\sigma}$. These are at most the required
source bounds. This handles the capacity-below-one case without
assuming it away.

Perform this construction on the \emph{ordered upper endpoints}
of each fine label vector at each required $U$, then freeze its
slot assignment for the whole box. The construction uses ordered
positive sizes, not primality of those endpoints. Original repeated
labels contribute their reciprocal factorial once. Each new group
convolution divided by its number-of-slots factorial is one-bounded;
restore the three factorials only in the atom's scalar. Attach beta
to group three. Its primes are disjoint from rough primes, so this
group is still one-bounded. There are finitely many atoms and a
bounded total absolute scalar mass independent of $T$.

The assignment may depend on the atom, $U$ and $T$ through endpoint
sizes. It does not depend on a particular prime tuple inside that
box or on the signs of the resulting progression errors. All these
are decompositions of one coefficient $\rho_T$, not replacements
of $\rho_T$ separately for different outer products.

\paragraph{5. The local-scale losses and the actual residue.}
Freeze the geometric margins and the least outer prime exponent
$a>0$. Choose $\zeta$ with $12\zeta<a/2$, then
$\epsilon\ll\zeta$ and $\kappa\ll\epsilon$. Work on
$x\ge T^{1-\kappa}/O(1)$. The medium slot has
$x^{2\zeta}\le U\le x^{1/3+\zeta/2}$, so
$A_0=Ux^{-\zeta}$ lies in
$[x^\zeta,x^{1/3-\zeta/2}]$. The endpoint construction gives the
three factorization monomials with savings $3\zeta,3\zeta,2\zeta$.
Replacing $A_0$ by $U$ leaves savings
\[
 \zeta,\quad 2\zeta,\quad(1-2\sigma)\zeta=9\zeta/16.
\]
Choose the beta-level and fine-endpoint losses smaller than these
fixed quantities, and $15\epsilon$ smaller than the remaining
minimum. Then all four inequalities
\eqref{eq:current8-source-sizes} hold. In particular the third
one retains $T^{7/32}$; it cannot be replaced by a constant.
All factor exponents in the fourth are positive, including
$5-2\sigma=73/16$, so smaller dyadic modulus lengths preserve it.
For $\kappa<\epsilon/(2(1+\epsilon))$ the actual residue $a=T$
satisfies $T<x^{1+\epsilon}$ eventually. This is uniform along
all even targets, not an average or a density-one subsequence.

\paragraph{6. What the outer coefficients actually are.}
Apply order-three HB separately at each original prime scale
$R_i\gg T^a$, retaining only dyadic rectangles meeting that
original group interval. Each group has at most six slots, and
its lower-endpoint product is at least $R_i/64$. The roughness
condition on the original product factors into the same condition
on every HB slot. A selected medium slot is therefore rough
$1$ or rough $\mu$, with an interval, normalized logarithm and
smooth Mellin twist. The full auxiliary SW lemma in
\texttt{focused-rough-sw-product-interface-audit.tex}, especially
its Steps1--3, proves precisely Definition2.1, uniformly for these
target-dependent geometric intervals and bounded variations.
It is not an SW assertion about an arbitrary grouped convolution.

Only the complementary divisor-bounded coefficient is truncated
at a logarithmic height and normalized. The discarded mass is
bounded \emph{after} the modulus sum by
\[
 (\log T)^{-B}\sum_{n\ll T,\ n\ne T}
       \tau_C(n)\tau_C(|T-n|)
       \ll T(\log T)^{C'-B}.
\]
Here $n$ is the whole HB product. Cauchy--Schwarz and the two
fixed divisor moments give the displayed uniform estimate.
The exact-coprime means have the analogous reciprocal-divisor
bound. For even $T$, every retained rough HB factor, including
units, is odd, so $n=T$ actually disappears termwise. The
squarefree-modulus collision error is paid before HB, on the
bounded positive outer tuple weight, by
$(T+T^{\theta_*})T^{-\omega}(\log T)^C$.

If there is no medium slot, every original group has a long slot;
otherwise its product is at most $x^{12\zeta+o(1)}$, too small.
There are at most two long slots. Thus this alternative has at most
two original prime factors, with residual product at most
$x^{20\zeta+o(1)}$. One long slot is handled by Lemma3.4, including
the exact primitive mean; two are handled by Proposition3.2, whose
residue supremum is over all nonzero integers. Its size condition
holds after choosing
$3/5<(1-20\zeta)(2/3-3\epsilon)$.
The required arbitrary individual subintervals occur explicitly
in Maynard II, arXiv:2006.07088v1, proof of Proposition5.2
(saved text lines1386--1467). Partial summation permits the two
smooth twists; it does not permit arbitrary arithmetic weights
on the long slots. Low rough modulus rank was separately handled
in paragraph2. Generalized BV below $1/2$ uses an original prime
SW factor and a complementary divisor-bounded factor; for one
original prime use classical BV.

This proves for the same compiled weight and each permitted smooth
geometric family $W_T$ the signed estimate
\[
 \sum_d\rho_T(d)\left\{
 \sum_{p_1\cdots p_k\equiv T\ (d)}W_T
 -\frac1{\varphi(d)}\sum_{(p_1\cdots p_k,d)=1}W_T\right\}
 =O_{A,W,\mathcal T}(T\log^{-A}T).
\]
Finite atom counts, their absolute scalars and logarithmically many
HB/dyadic/Mellin pieces are absorbed by asking for a larger saving.
The implied constants depend on the fixed policy and positive
cutoff, not on $T$. An arithmetic residue filter on an outer slot
would not follow from this proof. For example restricting a prime
slot to one residue class modulo three already destroys the
claimed unrestricted SW discrepancy. Divisors of $T$ are excluded
from the \emph{modulus} factors without this problem; if an outer
family needs them excluded too, first remove their bounded tuple
mass by the standard $O(T^{1-a+o(1)})$ estimate.

\paragraph{7. Positive mass, fixed choices and feedback.}
The pointwise upper polynomial is applied to a nonnegative global
majorant before the signed distribution estimate. For fixed $k$
and $p_i\ge T^a/2$,
\[
 \sum_{d:p\mid d}\frac{|\rho_T(d)|}{\varphi(d)}
 \ll\frac{\log T}{p-1}
\]
removes the exact coprime mass with error
$O_{a,k}(X_WT^{-a}\log T)$. Products below $T^{1-\kappa}$
cost $T^{1-\kappa+o(1)}$ after divisor summation and are not
assigned an unjustified local dispersion estimate. The radial
coordinate $(p_1\cdots p_k)/T$, as well as exponent coordinates,
is used in \texttt{prime-product-radial-boundary-repair.tex}; a
pure product wall is not null in the full exponent space. These
steps give $(KJ_W+o(1))C_TT/\log^2T$ from the actual density
bound $(K+o(1))C_T/\log T$. For the ordered bad-semiprime
family $J_W\to\log(1/b)$, with no extra factor two.

There is a simultaneous choice for the finite collection of current
feedback families. First reserve positive density/compiler losses
inside the retained endpoint slack. Choose a common compact trim
and coarse history partition, keeping all old decisions. Freeze
their strict gate margins. Next choose the dispersion parameters
above. Couple beta precision and rough cutoff, for example
$\ell=\nu_0\eta$ and $\omega=\nu_0\eta^2$, with rational
$\nu_0>0$. As $\eta\downarrow0$, both
$e^{-c/\eta}\omega^{-2}$ and the original third-tail bound
$O(\omega\log^2(1/\omega))$ tend to zero, while $\ell\to0$.
Thus the beta, tail and support losses fit the fixed budgets
simultaneously. Keep $\omega$ below every active marked cutoff.
Once the resulting rough-rank bound is fixed, choose a fine width
$\tau$ so that both the factorization loss $O(M\tau)$ and the
density error $C_{\mathcal T}\tau\omega^{-5}$ fit. This refines
coarse labels, rather than redefining them or changing a selected
gain. The safe total support loss can also be made below $1/500$.
Only then let $T\to\infty$.

The finite signed error is therefore absolute $o(\mathcal H_T)$,
where $\mathcal H_T=C_TT/\log^2T$ and $C_T$ is bounded below.
It is not a relative error required uniformly for arbitrarily
small cell mass. Every chosen cell/family and every finite number
of them have been fixed already. If endpoint integrals are used,
one chooses one sufficiently fine finite realization before this
limit; an $N$-adaptive mesh or depth is not licensed.

Finally \texttt{upper-k-to-current-lower-feedback-20261009.tex}
starts with the explicit hypothesis that the product theorem is
\emph{already proved} at the chosen fixed $K\le6.7618$. It does
not prove it for every smaller positive $K$. Once this hypothesis
is supplied by the upper construction, each of the five retained
positive cost occurrences can use the same theorem. Their fixed
switched bounds are affine in $K$, with the saved nonnegative
slopes; ordinary alternatives and old table values stay retained.
The table430 occurrence currently belongs only to $U_3,U_4$ and
is not part of that five-cost slope. The current change in $K$
is applied only after the upper theorem and its budgets have been
settled. This order creates no lower-to-upper circularity.

\paragraph{Completion boundary of this review.}
The selected-mask138 support-to-dispersion implication is supplied
above, including the small-third-prime rank split, target-dependent
modulus factors, the retained $T^{7/32}$, and absolute signed error.
No counterexample or unsatisfied new premise was found within
this class. The excluded stronger claims are: distribution for
arbitrary arithmetic outer weights; termwise-in-modulus estimates;
uniformity as $a\downarrow0$ or the policy refines with $T$;
automatic factorization at the carrier $\Theta$; and automatic
availability of every $K\le6.7618$. This review does not re-prove
external L23, ordinary raw-density convergence, the complete Chen
weight/table induction, or all87 numerical certificates. Those
are separate identified inputs, not consequences of this local
interface audit. No global constant or claim of literature priority
is changed.
